\section{总体思路与实现框架}

\begin{figure}[H]
  \centering
  \includegraphics[width=0.62\textwidth]{figures/fig_t4_directional_concept.pdf}
  \caption{定向源"方向不对"示意：波束迎光侧返回 \texttt{direction}，背光侧返回
  \texttt{no\_signal}，机器狗必须走到源正面才能听到。}
  \label{fig:concept}
\end{figure}

\subsection{检测阶段：以"半圆盘命中"为目标设计测量布局}

定向源在测量点 $p$ 处可被听到，当且仅当
\begin{equation}
  |p-g| \le R \quad\text{且}\quad (p-g)\cdot u(\theta) \ge 0,
  \qquad R\in[1000,1500]\,\text{m},\ \theta\in[0^\circ,360^\circ),
  \label{eq:dircond}
\end{equation}
其中 $u(\theta)=(\cos\theta,\sin\theta)$ 为波束方向单位向量。即检测区是半径 $R$ 的
\textbf{半圆盘} $H(g,\theta,R)$，而 $g,\theta,R$ 全部未知。检测布局的任务是使测量点集合
"尽量命中任意可能出现 $g,\theta,R$ 的半圆盘"。本文采用 \textbf{原点 + 7 覆盖基点 +
12 均匀方位外圈点} 的 20 点定案布局（第~\ref{sec:detect} 节），并给出 $422.5$ 万例
蒙特卡洛与对抗枚举的听到率统计。

\subsection{定位阶段：复用交会几何，剔除失效约束}

听到某频道后，其定位区域由 \texttt{direction}（楔形与接收半径内包）与 \texttt{near}
（$\le 5$ m）两类硬约束求交得到（$DirProbRegion$），\texttt{no\_signal} 一律不作为
约束。补测点按文献准则（示向度 Fisher 信息熵 / 交会几何度量，与问题三同源的
$t3.probing$ 薄转发）选取，多视角交会把"还差一点的"频道在清除途中收窄。

\subsection{清除阶段：顺路清除 + 就近试清 + 兜底逼近}

清除分三类：\textbf{(a) 途中顺路清除}——机器狗站间移动时，对"本站→下一站→圆心"夹出的
方位扇区内的未清频道估计点顺路试清（命中白捡时间，未中只花 $3$ s）；\textbf{(b) 就近
试清}——定位区域最小覆盖圆半径 $\le 400$ m 时直接试清，未中再用至多 $4$ 个半径 $20$ m
的圆覆盖定位区域逐点试清；\textbf{(c) 兜底逼近}——沿示向度逐步逼近直至 $\le 20$ m 清除
半径内。

\subsection{时间视角：耗时构成与下界验证}

虚拟耗时 $\approx$ 移动（$78\%$，其中检测扫描路线 $17.0$ km 约占 $54\%$）+ 测向
（$17\%$，大半是每位置必测未听频道的试探）+ 清除（$5\%$）。前两项本质不可压缩；
本文通过 20 局固定 seed 演练逐旋钮验证，收敛到 $6445$ s 的实测时间下界，并用神经网络
代理搜索独立交叉验证（第~\ref{sec:verify} 节）。